\chapter{Tenths and Hundredths}\label{ch:g4:decimals}

Between $3$ and $4$ the whole numbers are silent --- but a runner can
be $3$ meters and a bit ahead. Cutting the unit into ten, then a
hundred, gives the \emph{decimal fractions}, and a wonderful
shorthand for them: the decimal point. (Decimal numbers get their
full chapter in \cref{ch:g5:decimals}.)

\section{Decimal fractions}

\begin{definition}[Tenths and hundredths]\label{def:g4:decimals:def}
Cutting a unit into $10$ equal parts gives \emph{tenths},
$\frac{1}{10}$; cutting each tenth into $10$ again gives
\emph{hundredths}, $\frac{1}{100}$. So
\[
1 = \frac{10}{10} = \frac{100}{100},
\qquad
\frac{1}{10} = \frac{10}{100}.
\]
\end{definition}

\begin{omfigure}
\begin{tikzpicture}[scale=0.42]
  \draw[gray!50, thin] (0,0) grid (10,10);
  \foreach \i in {0,...,9} \fill[omDef!35] (3,\i+0.06)
    rectangle (3.94,\i+0.94);
  \fill[omThm!70] (6.05,4.05) rectangle (6.95,4.95);
  \draw[thick] (0,0) rectangle (10,10);
  \node[below, font=\small] at (5,-0.4)
    {the whole square $= 1$;\ \ a column (blue) $= \frac{1}{10}$;\ \
    a small cell (red) $= \frac{1}{100}$};
\end{tikzpicture}

{\small The hundred-square: ten columns of ten cells. One column is a
tenth, one cell a hundredth.}
\end{omfigure}

\begin{example}\label{ex:g4:decimals:count}
Shading $3$ columns and $7$ extra cells of the hundred-square shades
\[
\frac{3}{10} + \frac{7}{100} = \frac{30}{100} + \frac{7}{100}
= \frac{37}{100}
\]
of the square --- thirty-seven hundredths.
\end{example}

\section{The decimal point}

\begin{definition}[Decimal writing]\label{def:g4:decimals:point}
The number $2 + \frac{3}{10} + \frac{7}{100}$ is written
\[
2.37
\]
--- the \emph{decimal point} separates the whole part ($2$) from the
decimal part; the first digit after the point counts the tenths, the
second the hundredths. Reading: ``two point three seven'', or ``two
and thirty-seven hundredths''.
\end{definition}

\begin{example}\label{ex:g4:decimals:convert}
\[
\frac{7}{10} = 0.7, \qquad
\frac{43}{100} = 0.43, \qquad
\frac{5}{100} = 0.05, \qquad
3 + \frac{4}{10} = 3.4 .
\]
Careful with $\frac{5}{100}$: five \emph{hundredths} need a $0$ in
the tenths place --- $0.05$, not $0.5$.
\end{example}

\begin{omfigure}
\begin{tikzpicture}[scale=1.35]
  \draw[-Stealth] (-0.2,0) -- (8.4,0);
  \foreach \i in {0,10,20,30,40,50,60,70,80}
    \draw ({\i*0.1},-0.09) -- ({\i*0.1},0.09);
  \foreach \i in {0,...,80} \draw ({\i*0.1},-0.045) -- ({\i*0.1},0.045);
  \foreach \x/\l in {0/2, 1/{2.1}, 2/{2.2}, 3/{2.3}, 4/{2.4},
                     5/{2.5}, 6/{2.6}, 7/{2.7}, 8/{2.8}}
    \node[below, font=\footnotesize] at (\x,-0.1) {$\l$};
  \fill[omThm] (3.7,0) circle (1.7pt)
    node[above, font=\small] {$2.37$};
\end{tikzpicture}

{\small Zooming between $2$ and $2.8$: big marks every tenth, small
marks every hundredth. The number $2.37$ sits between $2.3$ and
$2.4$, at the $7$th small mark.}
\end{omfigure}

\begin{example}[Money is hundredths]\label{ex:g4:decimals:money}
One cent is one hundredth of one euro (or dollar, or pound): a price
of $4.85$ means $4$ whole coins and $85$ hundredths. Money is the
decimal point you already use every day: $10$ cents $= 0.10$; two
euros fifty $= 2.50$; and $0.05$ is a five-cent coin, not fifty!
\end{example}

\section{Comparing with tenths}

\begin{method}[Comparing decimals (first contact)]\label{met:g4:decimals:compare}
\begin{enumerate}
  \item Compare the whole parts first: $3.2 > 2.9$.
  \item Same whole parts: compare the tenths, then the hundredths:
        $2.37 < 2.5$ because $3$ tenths $< 5$ tenths.
  \item Writing both with two decimals helps: $2.5 = 2.50$, and
        $37 < 50$ hundredths.
\end{enumerate}
The trap: $2.37$ has more digits than $2.5$, but is
\emph{smaller} --- more digits does not mean bigger!
\end{method}

\begin{example}\label{ex:g4:decimals:compare}
Order $0.4$, $0.35$ and $0.09$: with hundredths, $0.40$, $0.35$,
$0.09$; so
\[
0.09 < 0.35 < 0.4 .
\]
\end{example}

\section{Exercises}

\begin{exercise}[$\star$]\label{exo:g4:decimals:1}
On a hundred-square, how much is shaded if one shades: $5$ columns;
$2$ columns and $3$ cells; $45$ cells? Answer with fractions.
\end{exercise}

\begin{exercise}[$\star$]\label{exo:g4:decimals:2}
Write with a decimal point:
$\frac{9}{10}$; \quad $\frac{27}{100}$; \quad $\frac{3}{100}$; \quad
$5 + \frac{6}{10}$; \quad $\frac{60}{100}$.
\end{exercise}

\begin{exercise}[$\star$]\label{exo:g4:decimals:3}
Write as a fraction (tenths or hundredths):
$0.3$; \quad $0.51$; \quad $0.07$; \quad $1.9$.
\end{exercise}

\begin{exercise}[$\star$]\label{exo:g4:decimals:4}
In $6.58$: what does the digit $5$ count? The digit $8$? Decompose
the number as in \cref{def:g4:decimals:point}.
\end{exercise}

\begin{exercise}[$\star$]\label{exo:g4:decimals:5}
Place on a number line graduated in tenths from $4$ to $5$:
$4.2$; \ $4.75$ (between which two marks?); \ $4.9$.
\end{exercise}

\begin{exercise}[$\star$]\label{exo:g4:decimals:6}
Copy and complete with $<$, $>$ or $=$:
\[
0.6 \;?\; 0.60, \qquad
2.8 \;?\; 2.75, \qquad
0.09 \;?\; 0.1, \qquad
5.3 \;?\; 5.13 .
\]
\end{exercise}

\begin{exercise}[$\star$]\label{exo:g4:decimals:7}
Order from smallest to biggest: $1.05$; \ $1.5$; \ $0.95$; \ $1.45$.
\end{exercise}

\begin{exercise}[$\star$]\label{exo:g4:decimals:8}
Write in decimal writing: three euros and five cents; twelve euros
and fifty cents; ninety-nine cents. Then order the three prices.
\end{exercise}

\begin{exercise}[$\star$]\label{exo:g4:decimals:9}
Which decimal numbers with one decimal digit lie strictly between
$3$ and $4$? List them all. How many are there?
\end{exercise}

\begin{exercise}[$\star\star$]\label{exo:g4:decimals:10}
Tom says: ``$0.25$ is bigger than $0.5$ because $25$ is bigger than
$5$''. Show him his mistake with the hundred-square (how many cells
does each shade?).
\end{exercise}

\begin{exercise}[$\star\star$]\label{exo:g4:decimals:11}
A runner's three jumps measured $3.6$ m, $3.58$ m and $3.65$ m.
Order the jumps from shortest to longest. Which two differ by
exactly $\frac{5}{100}$ of a meter?
\end{exercise}
